\section{Policy Gradient as Compute Allocation}

\smallskip

\begin{whythis}
You can derive the policy gradient theorem and implement REINFORCE. This chapter assumes that background. What it adds is the cost dimension: the same math that tells you \textit{which direction} to update the policy also tells you \textit{how much compute each design choice burns}. The PG literature optimizes reward. The job optimizes reward per GPU-hour --- those have different answers.

Which design choices actually move the cost needle, and by how much? Five choices dominate. This chapter derives each from the math, estimates its cost impact, and gives you an experiment on TinyRL (Chapter~7) to validate.
\end{whythis}

% -----------------------------------------------------------------------
\subsection{The RLVR Cost Ledger}
% -----------------------------------------------------------------------

One GRPO training step breaks down as:

\begin{lstlisting}
wall_clock_per_step =
    rollout_generation    (num_prompts * G * seq_len / throughput)    60-80%
  + reward_verification   (num_prompts * G * verifier_cost)           5-10%
  + reference_forward     (num_prompts * G * ref_model_fwd)           5-10%
  + backward_pass         (batch_size * model_size)                    5-15%
  + gradient_comm         (gradient_size / interconnect_bw)            2-5%
\end{lstlisting}

\begin{center}
\begin{tikzpicture}[x=0.065cm, y=0.8cm]
  % Stacked bar
  \fill[beni!25, rounded corners=1pt]   (0,0) rectangle (70,0.7);
  \fill[kitsune!20, rounded corners=1pt] (70,0) rectangle (80,0.7);
  \fill[kitsune!10!white, rounded corners=1pt] (80,0) rectangle (90,0.7);
  \fill[aiiro!20, rounded corners=1pt]   (90,0) rectangle (100,0.7);
  \fill[sumi!15, rounded corners=1pt]   (100,0) rectangle (105,0.7);

  % Labels above
  \node[font=\scriptsize, anchor=south] at (35, 0.75)  {Rollout generation (60--80\%)};
  \node[font=\scriptsize, anchor=south] at (85, 0.75)  {\tiny reward + ref};
  \node[font=\scriptsize, anchor=south] at (95, 0.75)  {\tiny bwd};

  % Brace below generation
  \draw[decorate, decoration={brace, amplitude=4pt, mirror}, thick, beni!50]
    (0, -0.15) -- (70, -0.15)
    node[midway, below=5pt, font=\scriptsize\itshape, beni!60] {where cost levers hit hardest};
\end{tikzpicture}
\end{center}

Ranges are regime-dependent --- they shift with model size, sequence length, and hardware. Chapter~5 covers the communication term in detail. The point is the ordering, not the exact numbers.

Rollout generation is 60--80\% of wall-clock. Cutting rollout count in half saves 30--40\% of total training cost. Cutting optimizer steps in half saves 3--8\%. The five levers in this chapter are ordered by this asymmetry.

\begin{thinkfirst}{Cost Asymmetry}
Two estimation problems. Work both on paper before reading the next section.

\textbf{(a)} Your GRPO pipeline runs $G = 16$ rollouts per prompt. You discover a variance reduction technique that gives equivalent gradient signal quality at $G = 8$. Using the cost ledger above, estimate the wall-clock savings as a percentage of total step time.

\textbf{(b)} You discover a second-order optimizer that converges in half the steps. Same question: estimate the wall-clock savings.

Which improvement would you staff an engineer on first?
\end{thinkfirst}

\medskip
\textbf{The five levers:}
\begin{enumerate}
  \item \textbf{Rollout allocation} --- how many rollouts per prompt, and which prompts get them
  \item \textbf{Regularization scheduling} --- when the KL penalty is miscalibrated for the training phase
  \item \textbf{Credit assignment} --- whether gradient signal targets the tokens that mattered
  \item \textbf{Preconditioning} --- whether the optimizer exploits policy-space curvature
  \item \textbf{Collapse detection} --- how early you detect that continued training is waste
\end{enumerate}

The first three are ordered by steady-state cost impact. The fourth targets the optimizer's share of wall-clock --- small but the only lever that reduces step count directly. The fifth prevents catastrophic waste rather than improving per-step efficiency --- a single early intervention can save more compute than all other levers combined.

% -----------------------------------------------------------------------
\subsection{Lever 1: Rollout Allocation}
% -----------------------------------------------------------------------

\begin{thinkfirst}{Optimal Group Size}
The GRPO gradient estimator for a single prompt with $G$ rollouts is:
\begin{align*}
  \hat{g} = \frac{1}{G} \sum_{i=1}^{G} \nabla_\theta \log \pi_\theta(\tau_i) \cdot \hat{A}_i, \quad \hat{A}_i = \frac{R_i - \bar{R}}{\sigma_R}
\end{align*}

The variance of this estimator scales as $\mathrm{Var}(\hat{g}) \propto \sigma^2_A / G$. Doubling $G$ halves variance but doubles the dominant cost term.

Derive: the cost-optimal $G$ minimizes total compute to reach a target reward. If convergence requires $N(G)$ steps and each step costs $c_{\mathrm{base}} + c_{\mathrm{rollout}} \cdot G$, the total cost is $N(G) \cdot (c_{\mathrm{base}} + c_{\mathrm{rollout}} \cdot G)$.

Under the simplifying assumption $N(G) \propto 1/G$ (linear variance reduction), cost is minimized as $G \to \infty$. Why does this not hold in practice? What determines the finite optimum?
\end{thinkfirst}

The assumption $N(G) \propto 1/G$ breaks because $N(G)$ flattens once gradient variance drops below the \textit{bias floor} --- the irreducible error from advantage estimation, stale rollouts, and reward noise. Below that floor, additional rollouts reduce variance that was not the binding constraint. The real optimum is finite and depends on the prompt distribution.

\textbf{Waste pattern: dead rollouts.} GRPO normalizes advantages within each group: $\hat{A}_i = (R_i - \bar{R}) / \sigma_R$. For prompts where all $G$ rollouts get similar rewards ($\sigma_R \approx 0$), normalized advantages are near-zero. Those rollouts produce negligible gradient signal. If 30\% of prompts have low reward variance --- easy problems the model already solves reliably, or hard problems it consistently fails --- then $\sim$30\% of the rollout budget produces near-zero learning.

\textbf{Cost impact.} Napkin math: a 7B model GRPO run at $G = 16$ with 100K prompts over 5K steps generates 8 billion rollout tokens. If 30\% are dead rollouts, that is 2.4 billion tokens of inference for near-zero gradient signal. At $\sim$\$0.01 per 1K inference tokens, that is $\sim$\$24K wasted on a single run.

\textbf{Systems reality.} $G$ rollouts from the same prompt share the prompt's KV cache --- the prompt forward pass is computed once and reused across all $G$ completions. Going from $G = 8$ to $G = 16$ adds 8 completion forward passes, not 8 full forward passes. This makes larger $G$ cheaper than the naive cost model predicts.

But: KV cache for $G$ completions must fit in HBM simultaneously. At $G = 16$ with long prompts, you can hit memory limits that force smaller batch sizes, which hurts GPU utilization. The real constraint is $G \times \text{batch\_size} \times \text{seq\_len} \leq \text{HBM budget}$.

\textbf{Literature (established solutions).} This problem is solved. Multiple papers from early 2026 address adaptive rollout allocation:

\begin{itemize}
  \item \textbf{VIP} (\href{https://arxiv.org/abs/2602.01601}{arXiv:2602.01601}, 2026): Variance-Informed Predictive allocation. A Gaussian process predicts per-prompt success probabilities from recent rollouts. These predictions feed a convex optimization problem that allocates rollout budget to minimize expected gradient variance under a hard compute constraint.
  \item \textbf{AERO} (\href{https://arxiv.org/abs/2602.14338}{arXiv:2602.14338}, 2026): Adaptive rollout strategy with selective rejection and a Bayesian posterior to prevent zero-advantage dead zones. Reports 48\% compute reduction on 1.5B and 47\% on 7B models.
  \item \textbf{Prompt Replay} (\href{https://arxiv.org/abs/2603.21177}{arXiv:2603.21177}, 2026): On-policy reuse of high-signal prompts --- skip prompts that have become too easy or too hard, replay prompts in the high-learning zone.
\end{itemize}

\begin{exercise}{Rollout Allocation Efficiency (4--6 hours)}

\textbf{Part 1: Fixed $G$ sweep.} Train GRPO on TinyRL-Code (Chapter~7, 124M config) with $G \in \{4, 8, 16, 32\}$. For each $G$, measure:
\begin{itemize}
  \item Reward per GPU-second (primary metric)
  \item Dead rollout rate: fraction of prompts per step where $\sigma_R < 0.1$
  \item Gradient norm variance across steps
\end{itemize}
Plot reward/GPU-second vs.\ $G$. Is there a clear optimum?

\textbf{Part 2: Adaptive allocation.} Implement a simplified VIP-style scheme: maintain a rolling window of per-prompt reward variance. Each step, allocate $G_i \propto \hat{\sigma}_{R,i}$ under a fixed total rollout budget equal to $G = 8 \times \text{num\_prompts}$. Compare reward/GPU-second against the best fixed $G$ from Part 1.

\textbf{Part 3: Analysis.} Does the dead rollout rate increase over training (as the policy improves on easy prompts)? If so, a fixed $G$ becomes increasingly wasteful. Plot dead rollout rate over training steps for each fixed $G$. Does adaptive allocation keep the dead rollout rate low?

\textbf{What to investigate next.} Read VIP and AERO in full. Both assume reward variance is a good proxy for gradient information --- is this true when the advantage estimator introduces its own bias? Compare VIP's Gaussian process approach against AERO's Bayesian posterior. Which is cheaper to maintain at scale?
\end{exercise}

% -----------------------------------------------------------------------
\subsection{Lever 2: Regularization Scheduling}
% -----------------------------------------------------------------------

\begin{thinkfirst}{The KL Penalty as Resource Allocation}
GRPO adds a KL penalty to the reward:
\begin{align*}
  R'(\tau) = R(\tau) - \beta \cdot D_{\mathrm{KL}}(\pi_\theta \| \pi_{\mathrm{ref}})
\end{align*}
This is equivalent to optimizing a variational lower bound:
\begin{align*}
  \max_\theta \; \mathbb{E}_{\pi_\theta}[R(\tau)] - \beta \cdot D_{\mathrm{KL}}(\pi_\theta \| \pi_{\mathrm{ref}})
\end{align*}

The coefficient $\beta$ controls how far the policy drifts from the reference. A fixed $\beta$ applies the same exploration-stability trade-off for the entire run.

Think through the consequences at two extremes of training:
\begin{enumerate}
  \item \textbf{Early training:} the policy is far from any reward-hacking regime and the gradient signal is noisy. Does a fixed $\beta$ help or hurt?
  \item \textbf{Late training:} the policy is near a good solution. Small perturbations can trigger entropy collapse. Does the same $\beta$ help or hurt?
\end{enumerate}
What would an adaptive schedule look like?
\end{thinkfirst}

\textbf{Waste pattern.} Over-regularization wastes gradually: each step is slightly less efficient because $\beta$ constrains updates that would have been safe. Under-regularization wastes catastrophically: the policy drifts into collapse or reward-hacking, and hundreds of steps are spent recovering. The cost of getting $\beta$ wrong is dominated by the catastrophic tail.

Early in training, $\beta$ barely constrains --- the reference forward pass (5--10\% of step time per the cost ledger) produces no useful signal. Late in training, the same $\beta$ may be too loose, causing collapse-and-recovery episodes that consume hundreds of steps.

\textbf{Systems reality.} The reference model consumes HBM. Common workarounds and their cost profiles:
\begin{itemize}
  \item \textbf{CPU offload:} frees HBM but adds latency for each KL computation.
  \item \textbf{Weight sharing with LoRA:} reference = frozen base, actor = base + adapter. Eliminates the duplicate model but constrains the policy to a low-rank offset.
  \item \textbf{Asynchronous KL:} compute KL on a separate device or with a stale snapshot. Reduces latency on the critical path but introduces approximation error.
\end{itemize}
Each workaround changes the cost of evaluating KL, which changes when adaptive scheduling breaks even.

\textbf{Literature:}
\begin{itemize}
  \item ``Rethinking KL Regularization in RLHF'' (\href{https://arxiv.org/abs/2510.01555}{arXiv:2510.01555}, 2025): argues GRPO's KL implementation overlooks its functional role as optimization loss.
  \item \textbf{SAFE} (\href{https://arxiv.org/abs/2602.04651}{arXiv:2602.04651}, 2026): PID-controlled adaptive KL thresholds with entropy-gated regulation. Formalizes adaptive $\beta$ using control theory.
  \item \textbf{APO} (\href{https://arxiv.org/abs/2512.22953}{arXiv:2512.22953}, 2025): $\alpha$-divergence interpolation between forward and reverse KL with confidence-guarded scheduling.
\end{itemize}

\begin{exercise}{KL Scheduling Comparison (4--5 hours)}

Train GRPO on TinyRL-Code (124M) under four conditions:
\begin{enumerate}[label=(\alph*)]
  \item Fixed $\beta = 0.01$
  \item Fixed $\beta = 0.1$
  \item Simple adaptive: multiply $\beta$ by 1.5 when $D_{\mathrm{KL}} > 0.05$, by 0.67 when $D_{\mathrm{KL}} < 0.01$
  \item PID-controlled $\beta$ (reproduce SAFE's approach): target $D_{\mathrm{KL}} = 0.03$, proportional gain $K_p = 1.0$, integral gain $K_i = 0.1$, derivative gain $K_d = 0.01$
\end{enumerate}

Measure at every step: reward/GPU-second, $D_{\mathrm{KL}}(\pi_\theta \| \pi_{\mathrm{ref}})$, policy entropy. Count entropy collapse events (entropy $< 0.1$ for $> 10$ consecutive steps).

\textbf{Questions.} Does PID outperform the simple heuristic? By enough to justify the implementation complexity? Does adaptive scheduling prevent collapse episodes that the fixed conditions hit? Which fixed $\beta$ is closer to the adaptive performance?

\textbf{What to investigate next.} Read SAFE and APO. SAFE uses PID control --- is the integral term necessary, or does proportional control suffice? APO schedules across divergence families, not just $\beta$ magnitude --- does the choice of forward vs.\ reverse KL matter more than the coefficient?
\end{exercise}

% -----------------------------------------------------------------------
\subsection{Lever 3: Credit Assignment}
% -----------------------------------------------------------------------

\begin{thinkfirst}{Where Does the Gradient Signal Go?}
With sequence-level rewards, the policy gradient applies the same advantage $\hat{A}$ to every token:
\begin{align*}
  \hat{g} = \sum_{t=1}^{T} \nabla_\theta \log \pi_\theta(a_t | s_t) \cdot \hat{A}
\end{align*}

In a 512-token math reasoning trace, suppose $\sim$20 tokens are critical (an algebraic manipulation, a key substitution). The other 492 tokens are setup or mechanical follow-through.

\textbf{(a)} What fraction of the gradient variance comes from non-critical tokens? (Each token contributes $\|\nabla_\theta \log \pi_\theta(a_t | s_t)\|^2 \cdot \hat{A}^2$ to the variance.)

\textbf{(b)} How does this connect to the rollout budget from Lever~1? If gradient variance is 25$\times$ higher than it needs to be, how many extra rollouts does that require to compensate?
\end{thinkfirst}

The naive story --- ``GRPO only does sequence-level credit'' --- is wrong. Recent work shows GRPO already performs implicit token-level credit assignment.

\textbf{The implicit credit mechanism.} Sullivan et al.\ (\href{https://arxiv.org/abs/2509.21154}{arXiv:2509.21154}, ``GRPO is Secretly a PRM,'' 2025) prove that GRPO with outcome rewards is equivalent to a process-reward-model-aware RL objective with Monte-Carlo-based process rewards, whenever trajectories within a group share prefixes. In practice, $G$ rollouts from the same prompt often share long prefixes --- the model generates similar openings. Where prefixes overlap, GRPO assigns differential credit to the divergence points: implicit token-level credit without a critic.

This implicit mechanism breaks down when:
\begin{itemize}
  \item Prefixes are short (the model generates diverse openings --- common for creative or open-ended tasks)
  \item $G$ is small (fewer shared prefixes)
  \item The critical decision is early in the sequence (before prefixes diverge)
\end{itemize}

\textbf{Waste pattern.} When the implicit mechanism fails, you pay the full gradient variance tax from uniform credit. This inflates the required $G$ for stable training --- linking back to Lever~1. Poor credit assignment is a hidden multiplier on rollout cost.

\textbf{Explicit alternatives.} Multiple methods improve on implicit credit, each with a different compute profile:
\begin{itemize}
  \item \textbf{GRPO-$\lambda$} (OpenReview): eligibility traces using token-level log-probabilities. Explicit credit without a learned critic. Negligible overhead --- reuses existing log-probs.
  \item \textbf{GTPO/GRPO-S} (\href{https://arxiv.org/abs/2508.04349}{arXiv:2508.04349}, 2025): entropy-weighted per-token reward shaping via Dynamic Entropy Weighting. Small overhead (entropy computation per token).
  \item \textbf{SPO} (\href{https://arxiv.org/abs/2505.23564}{arXiv:2505.23564}, 2025): segment-level advantages. More precise than sequence-level, fewer estimation points than token-level. Moderate overhead (segment boundary detection).
  \item \textbf{Learned critic} (PPO-style): $\sim$30--50\% memory overhead $+$ $\sim$20\% step time for critic forward and backward passes.
\end{itemize}

The trade-off: explicit credit methods add per-step cost but reduce required $G$. Whether this nets positive depends on the ratio of rollout cost to credit-method cost.

\begin{exercise}{Credit Assignment Strategies (5--6 hours)}

Train GRPO on TinyRL-Code (124M) with four credit assignment strategies:
\begin{enumerate}[label=(\alph*)]
  \item Standard GRPO (implicit credit via shared prefixes)
  \item GRPO-$\lambda$ (eligibility traces)
  \item GTPO (entropy-weighted token rewards)
  \item Heuristic: assign zero advantage to tokens outside code blocks (task-structure proxy for ``where the reasoning happens'')
\end{enumerate}

Measure: reward/GPU-second (including all overhead), gradient norm variance, convergence speed to 80\% solve rate.

\textbf{Key question.} At what $G$ does standard GRPO (a) match GRPO-$\lambda$'s (b) performance at $G = 8$? The gap in $G$ is the credit assignment tax in rollout budget.

\textbf{Prefix analysis.} Measure the actual prefix-sharing rate in your runs: what fraction of rollout pairs within a group share $> 50$\% of tokens? If high, implicit credit may be sufficient and explicit methods add cost for little gain. If low, GRPO-$\lambda$ or SPO become high-value.

\textbf{What to investigate next.} Read ``GRPO is Secretly a PRM'' carefully. The proof relies on shared prefixes. Does prefix-sharing rate change over training (as the policy becomes more deterministic)? If so, implicit credit improves as training progresses --- exactly when it matters most.
\end{exercise}

% -----------------------------------------------------------------------
\subsection{Lever 4: Preconditioning}
% -----------------------------------------------------------------------

The weakest lever by cost impact. No published results apply second-order optimizers to policy gradient updates as of April 2026.

\begin{thinkfirst}{Natural Gradient and Distribution Geometry}
The ordinary policy gradient $\nabla_\theta J$ treats all parameter directions equally. But a 0.01 change in one parameter might shift the policy distribution dramatically (a weight gating a high-probability token), while 0.01 in another direction does nothing (a weight in a saturated layer).

The Fisher information matrix measures this sensitivity:
\begin{align*}
  F(\theta) = \mathbb{E}_{a \sim \pi_\theta(\cdot|s)}\!\left[\nabla_\theta \log \pi_\theta(a|s) \cdot \nabla_\theta \log \pi_\theta(a|s)^\top\right]
\end{align*}

The natural gradient $\tilde{g} = F(\theta)^{-1} \nabla_\theta J$ rescales by local curvature --- equal-sized steps in distribution space regardless of parameterization.

\textbf{(a)} PPO clips the importance ratio uniformly: the same $\epsilon$ for every token. The Fisher defines a trust region in KL space: $\delta^\top F \delta \leq c$. Why are these not equivalent? What does PPO miss?

\textbf{(b)} Adam approximates diagonal curvature --- it adapts per-parameter learning rates but ignores cross-parameter structure. In high-dimensional policy networks, the Fisher has low-rank structure: most curvature is concentrated in a small subspace. What does this imply about how much of Adam's adaptation budget is wasted?
\end{thinkfirst}

\textbf{Waste pattern.} Each optimizer step moves less in distribution space than it could. The step count to reach a target policy change is inflated. Per-step cost is unchanged.

\textbf{Cost impact.} Backward $+$ optimizer is 5--15\% of wall-clock (per the cost ledger). Even halving step count saves 2.5--7.5\% of total. This is the fourth lever for a reason.

\textbf{Literature (pretraining, not PG).}
\begin{itemize}
  \item \textbf{SOAP} (\href{https://arxiv.org/abs/2409.11321}{arXiv:2409.11321}, 2024): Shampoo $+$ Adam in preconditioner eigenbasis. 40\% iteration reduction, 35\% training time reduction for LLM pretraining.
  \item \textbf{KL-Shampoo} (\href{https://arxiv.org/abs/2509.03378}{arXiv:2509.03378}, 2025): outperforms SOAP on pretraining benchmarks.
  \item No published results on PG/RL applications as of April 2026.
\end{itemize}

\textbf{Systems reality.} Shampoo stores preconditioner matrices of shape $[d_{\mathrm{in}}, d_{\mathrm{in}}]$ and $[d_{\mathrm{out}}, d_{\mathrm{out}}]$ per layer. Memory overhead proportional to $d_{\mathrm{in}}^2 + d_{\mathrm{out}}^2$ --- significant for large models. Preconditioner updates require matrix root computations. In distributed training, preconditioner statistics must be aggregated across workers (additional All-Reduce cost; see Chapter~5).

\begin{exercise}{Preconditioning for Policy Gradients (4--6 hours)}

This is exploratory --- there are no published baselines for comparison.

Train GRPO on TinyRL-Code with three optimizer configurations:
\begin{enumerate}[label=(\alph*)]
  \item AdamW ($\beta_1 = 0.9$, $\beta_2 = 0.999$)
  \item AdamW with grid-searched LR schedule ($3$ learning rates $\times$ $3$ warmup lengths --- 9 short runs, pick the best)
  \item SOAP (if available in your framework) or block-diagonal K-FAC approximation
\end{enumerate}

Measure: reward/GPU-second, steps to 80\% solve rate, peak HBM usage. Run on both 124M and 350M configs.

\textbf{Key question.} Does SOAP's step-count reduction outweigh its per-step overhead? Does the answer change with model size?

\textbf{What to investigate next.} The comparison between well-tuned AdamW (b) and SOAP (c) is the decisive test. If a good LR schedule closes the gap, preconditioning is not worth the complexity for PG. If SOAP wins despite overhead, the next question is whether the gain scales with model size.
\end{exercise}

% -----------------------------------------------------------------------
\subsection{Lever 5: Collapse Detection}
% -----------------------------------------------------------------------

Entropy collapse is well-studied with specific mechanistic findings. This section teaches the established results; the experiment validates rather than discovers.

\begin{thinkfirst}{Why Does Clipping Cause Entropy Collapse?}
Entropy collapse occurs when $H(\pi_\theta(\cdot|s)) \to 0$: the policy becomes near-deterministic. When this happens, $\nabla_\theta \log \pi_\theta(a|s)$ diverges for low-probability actions, creating extreme gradient variance. The policy is stuck.

The cause is structural, not just a symptom of reward overfitting. Consider PPO/GRPO clipping with the importance ratio $\rho_t = \pi_\theta(a_t|s_t) / \pi_{\mathrm{old}}(a_t|s_t)$:

\begin{itemize}
  \item \textbf{Clip-high} (capping $\rho_t$ for negative advantages, $\hat{A}_t < 0$): prevents the policy from \textit{decreasing} probabilities too fast. But this asymmetrically allows probabilities to \textit{increase} without limit for positive advantages, concentrating probability mass. Net effect: entropy decreases.
  \item \textbf{Clip-low} (capping $\rho_t$ for positive advantages, $\hat{A}_t > 0$): prevents probabilities from \textit{increasing} too fast. This limits concentration. Net effect: entropy increases.
\end{itemize}

Under standard symmetric clipping ($\epsilon = 0.2$), the clip-high entropy-reducing effect dominates. Entropy decreases even with random rewards (Engstrom et al., ``Clip-Low Increases Entropy, Clip-High Decreases Entropy,'' 2025).

\begin{center}
\begin{tikzpicture}[
  x=0.8cm, y=2.5cm,
  every node/.style={font=\scriptsize}
]
  % Axes
  \draw[-{Stealth[length=4pt]}, sumi!50] (-0.3,0) -- (7,0) node[right] {tokens};
  \draw[-{Stealth[length=4pt]}, sumi!50] (0,-0.05) -- (0,1.2) node[above] {$\pi(a|s)$};

  % Before clipping — spread distribution
  \draw[aiiro!60, thick, smooth] plot coordinates
    {(0.5,0.05) (1,0.15) (2,0.35) (3,0.55) (4,0.35) (5,0.15) (6,0.05)};
  \node[aiiro!60, above] at (3, 0.58) {before};

  % After clip-high only — concentrated
  \draw[beni!60, thick, smooth, dashed] plot coordinates
    {(0.5,0.02) (1,0.05) (2,0.15) (3,0.85) (4,0.15) (5,0.03) (6,0.01)};
  \node[beni!60, above] at (4.5, 0.35) {after clip-high};

  % Annotation
  \draw[-{Stealth[length=3pt]}, beni!40, thick] (2.5, 0.3) -- (2.9, 0.6);
  \node[beni!40, anchor=east, text width=2cm, align=right] at (2.3, 0.45)
    {mass\\concentrates};
\end{tikzpicture}
\end{center}

\smallskip
\noindent\textit{Symmetric clipping's net effect: clip-high allows probability increases without bound for positive advantages, concentrating mass. Entropy decreases.}

\textbf{Work through:} if clip-high is the root cause, what does asymmetric clipping ($\epsilon_{\mathrm{low}} = 0.2$, $\epsilon_{\mathrm{high}} = 0.1$) do to the entropy trajectory?
\end{thinkfirst}

Token-level analysis (\href{https://arxiv.org/abs/2510.10150}{arXiv:2510.10150}, ``Rethinking Entropy Interventions,'' 2025): entropy change per update is governed by four factors --- clipping strategy, advantage sign, token probability, and token entropy. This gives a precise diagnostic: you can predict which tokens will lose entropy fastest.

\textbf{Waste pattern.} Every step after collapse begins is wasted compute. In long runs, collapse might begin at step 3000 but not be noticed until step 5000. Recovery is difficult: Adam's moment estimates have adapted to the collapsed regime, carrying stale curvature information.

\textbf{Cost impact.} Early detection by 500 steps saves $500 \times$ full step cost. At scale, this can exceed the savings from all other levers combined.

\textbf{Systems reality.} Detection must be cheap and non-interrupting:
\begin{itemize}
  \item Policy entropy over the training batch: free (logits already computed).
  \item Effective rank of logit matrix (via SVD): expensive for large vocab. Log every 50--100 steps to amortize.
  \item Reward std across rollouts ($\sigma_R \to 0$ means identical outputs): cheap, already available.
  \item Top-1 token probability (mean over batch): cheap proxy for determinism.
\end{itemize}

Intervention options with different cost profiles:
\begin{itemize}
  \item \textbf{Increase entropy bonus:} free. May not work if collapse is deep.
  \item \textbf{Revert to checkpoint $N$ steps before onset:} wastes $N$ steps but guarantees recovery.
  \item \textbf{Asymmetric clipping:} increase $\epsilon_{\mathrm{low}}$, decrease $\epsilon_{\mathrm{high}}$. Free, addresses root cause per the mechanism literature.
  \item \textbf{Partially reset Adam moments} for output layer parameters: moderate cost, targeted.
\end{itemize}

\textbf{Literature:}
\begin{itemize}
  \item ``The Entropy Mechanism of RL for Reasoning LMs'' (\href{https://arxiv.org/abs/2505.22617}{arXiv:2505.22617}, PRIME-RL, 2025): sharp early entropy drop leads to performance saturation.
  \item ``Clip-Low Increases Entropy, Clip-High Decreases Entropy'' (OpenReview, 2025): clipping mechanism itself causes entropy reduction.
  \item ``Rethinking Entropy Interventions in RLVR'' (\href{https://arxiv.org/abs/2510.10150}{arXiv:2510.10150}, 2025): four-factor token-level entropy change model.
  \item ``Flexible Entropy Control in RLVR'' (\href{https://arxiv.org/abs/2602.09782}{arXiv:2602.09782}, 2026): gradient-preserving entropy control methods.
\end{itemize}

\begin{exercise}{Validating the Clipping-Entropy Mechanism (4--5 hours)}

\textbf{Part 1: Reproduce the finding.} Train GRPO on TinyRL-Code (124M) with standard symmetric clipping ($\epsilon = 0.2$). Log at every step: policy entropy (mean over batch), reward std across rollouts, top-1 token probability, gradient norm. Log effective rank every 50 steps.

\textbf{Part 2: Asymmetric clipping.} Train with $\epsilon_{\mathrm{low}} = 0.2$, $\epsilon_{\mathrm{high}} = 0.1$. Compare the entropy trajectory against Part~1. Does asymmetric clipping slow or prevent entropy collapse? What happens to final reward --- does preserving entropy help or hurt?

\textbf{Part 3: Automated detection and intervention.} Implement threshold-based detection: trigger when entropy drops below its 200-step moving average by $> 2$ standard deviations. Compare two interventions:
\begin{enumerate}[label=(\roman*)]
  \item Switch to asymmetric clipping ($\epsilon_{\mathrm{low}} = 0.2$, $\epsilon_{\mathrm{high}} = 0.1$)
  \item Revert to checkpoint 200 steps before the trigger
\end{enumerate}
Which recovers faster? Which reaches higher final reward?

\textbf{Analysis.} Which metric first signals collapse onset? How many steps of warning does the earliest signal give before entropy drops below 0.1?

\textbf{What to investigate next.} The four-factor model from ``Rethinking Entropy Interventions'' predicts which tokens lose entropy fastest. Can you build a targeted intervention --- adjust clipping only for high-risk tokens rather than globally?
\end{exercise}

% -----------------------------------------------------------------------
\subsection{Measurement \& Reliability}
% -----------------------------------------------------------------------

The five levers are not independent knobs. Rollout allocation interacts with credit assignment (poor credit inflates required $G$). Regularization scheduling interacts with collapse detection (too-loose $\beta$ triggers the failure that Lever~5 catches). This section ties them together through what an RE monitors during a training run.

\textbf{The RLVR dashboard.} Minimum metrics to log at every step:

\begin{table}[h]
\centering
\small
\begin{tabular}{@{}lll@{}}
\toprule
\textbf{Metric} & \textbf{What it tells you} & \textbf{Lever} \\
\midrule
Reward / GPU-second    & Overall cost efficiency        & All \\
Mean reward (train)    & Learning progress              & --- \\
Mean reward (held-out) & Generalization / reward hacking & 2, 5 \\
Reward std per prompt  & Dead rollout rate              & 1 \\
Policy entropy (mean)  & Collapse proximity             & 5 \\
KL from reference      & Regularization regime          & 2 \\
Gradient norm          & Training stability             & 4, 5 \\
Top-1 token prob       & Determinism proxy              & 5 \\
\bottomrule
\end{tabular}
\end{table}

\textbf{Failure modes beyond entropy collapse.}

\begin{itemize}
  \item \textbf{Reward over-optimization.} Held-out reward diverges from training reward. The policy is specializing to the training prompt distribution or exploiting verifier-specific patterns. \textit{Signal:} training reward rises while held-out reward plateaus or drops. \textit{Intervention:} increase $\beta$, diversify prompt set.

  \item \textbf{Length gaming.} Completions grow longer without improving quality --- reward per token drops. The policy discovers that longer outputs score slightly higher due to verifier artifacts. \textit{Signal:} mean completion length increases while reward/token decreases. \textit{Intervention:} add length penalty to reward, or normalize reward by completion length.

  \item \textbf{Gradient explosion / NaN.} Triggered by rare high-reward outliers with extreme log-probabilities. A single rollout with reward $10\times$ the batch mean creates a gradient spike that destabilizes training. \textit{Signal:} gradient norm spikes by $> 10\times$ its moving average. \textit{Intervention:} gradient clipping (already standard), reward clipping, or outlier filtering in the advantage computation.

  \item \textbf{Straggler amplification.} In distributed rollout generation, one slow completion (e.g., a prompt that elicits a very long response) holds up the entire batch. With $G = 16$ and 100 prompts, one 4096-token completion among 1600 that average 512 tokens wastes GPU-seconds across all workers. \textit{Signal:} high variance in per-prompt generation time. \textit{Intervention:} generation length caps, asynchronous rollout collection, or padding-aware batching. See Chapter~5 for the distributed systems perspective.
\end{itemize}

\textbf{Logging cost as compute allocation.} Not everything should be logged every step. Effective rank requires SVD --- expensive for large vocabularies. Log every 50--100 steps. Held-out reward requires inference on a separate prompt set --- log every 200--500 steps. Dashboard design is itself a cost decision.

% -----------------------------------------------------------------------
\subsection{Interview Arsenal}
% -----------------------------------------------------------------------

\begin{taste}
The goal of this chapter is not to memorize five levers. It is to internalize a way of thinking: every algorithmic choice is a compute allocation decision, and the right way to evaluate it is to measure its cost impact, not its theoretical elegance.

The three interview responses below are not scripts to recite. They are structures to fill with your own experimental results from TinyRL. An interviewer who hears ``I ran this experiment and here is what I found'' will react differently from one who hears ``I read that this should work.''
\end{taste}

\textbf{``What would you investigate in your first two weeks?''}

\begin{enumerate}
  \item \textit{Profile the pipeline:} measure the actual time split across cost ledger terms on the specific hardware and model. Do not assume the 60--80\% rollout figure --- measure it.
  \item \textit{Measure the dead rollout rate:} what fraction of prompts have near-zero reward variance? This gives the ceiling for Lever~1 (rollout allocation).
  \item \textit{Check the KL trajectory:} is $\beta$ well-calibrated? Is there a visible collapse-and-recovery pattern in the training logs?
  \item \textit{Compute reward/GPU-second} across recent checkpoints. Improving, flat, or declining? Declining means one of the five waste patterns is active.
\end{enumerate}

\textbf{``What do you think policy gradient is really doing?''}

The answer should demonstrate three things:
\begin{enumerate}
  \item That PG is stochastic optimization under non-stationarity --- the sampling distribution changes with the parameters, unlike supervised learning where the data distribution is fixed.
  \item That each standard technique (baselines, clipping, KL penalties, natural gradient) addresses a specific consequence of that non-stationarity.
  \item A concrete opinion on which technique has the highest return on compute, backed by your experimental evidence from TinyRL or by specific literature results.
\end{enumerate}

\textbf{``How would you reduce training cost by 30\%?''}

Walk through the levers in cost-impact order with specific measurements. The structure: measure the pipeline, identify the largest waste source, run a controlled experiment before committing to a production change. The answer is never ``use a better algorithm'' --- it is ``here is the waste pattern I identified, here is the evidence, and here is the experiment I would run to validate the fix.''

\newpage
